% Canonical Paper II source.
We now specialize the surface theory to the complete model imported from
Paper~I.  The calculation is elementary, but fixing every scale at this point is
essential: the same constants enter the Poisson kernel, the conormal derivative,
and the boundary energy in the next two sections.

Throughout this section and Sections~\ref{sec:poisson-dirichlet}--
\ref{sec:boundary-energy}, assume
\[
 \mu>0,\qquad \lambda>0,
\]
and let
\[
 \HH=\{(x,y)\in\R^2:y>0\},\qquad
 g_{\mu,\lambda}
 =\frac1{y^2}\left(\frac{dx^2}{\mu^2}
                    +\frac{dy^2}{\lambda^2}\right),
 \qquad
 a(x,y)=-\frac{x}{y}.
\]
The orientation is that of $dx\wedge dy$.  We continue to use the
divergence-of-gradient convention for $\Delta_g$.

\subsection{Conformal normalization and the arithmetic frame}

Define the globally scaled coordinate and its complexification by
\begin{equation}\label{eq:hyperbolic-scaled-coordinate}
 X=\frac{\lambda}{\mu}x,
 \qquad
 w=X+iy.
\end{equation}

\begin{proposition}[Analytic normalization of the basic model]
\label{prop:hyperbolic-analytic-normalization}
In the coordinates \eqref{eq:hyperbolic-scaled-coordinate},
\begin{align}
 g_{\mu,\lambda}
   &=\frac1{\lambda^2}\frac{dX^2+dy^2}{y^2},
   \label{eq:hyperbolic-conformal-metric}\\
 d\vol_g
   &=\frac{dx\,dy}{\mu\lambda y^2}
    =\frac{dX\,dy}{\lambda^2y^2},
   \label{eq:hyperbolic-volume}\\
 \Delta_g
   &=\mu^2y^2\partial_x^2+\lambda^2y^2\partial_y^2
    =\lambda^2y^2(\partial_X^2+\partial_y^2).
   \label{eq:hyperbolic-laplace-beltrami}
\end{align}
For $F\in C^1(\HH;\C)$, its Dirichlet energy, without a factor $1/2$, is
\begin{equation}\label{eq:hyperbolic-dirichlet-energy}
\begin{aligned}
 \mathcal E_g(F)
 :=\int_{\HH}|dF|_g^2\,d\vol_g
 &=\int_{\HH}
   \left(\frac{\mu}{\lambda}|F_x|^2
        +\frac{\lambda}{\mu}|F_y|^2\right)dx\,dy
 \\
 &=\int_{\HH}(|F_X|^2+|F_y|^2)\,dX\,dy,
\end{aligned}
\end{equation}
whenever any one of these integrals is finite.
\end{proposition}

\begin{proof}
Since $dx=(\mu/\lambda)dX$, substitution into the metric gives
\eqref{eq:hyperbolic-conformal-metric}.  In the original coordinates,
\[
 (g_{ij})=
 \begin{pmatrix}
  (\mu^2y^2)^{-1}&0\\
  0&(\lambda^2y^2)^{-1}
 \end{pmatrix},
 \qquad
 \sqrt{\det g}=\frac1{\mu\lambda y^2}.
\]
This proves \eqref{eq:hyperbolic-volume}.  Moreover,
\[
 \sqrt{\det g}\,g^{xx}=\frac{\mu}{\lambda},
 \qquad
 \sqrt{\det g}\,g^{yy}=\frac{\lambda}{\mu}
\]
are constant.  The coordinate formula
\[
 \Delta_gF
 =\frac1{\sqrt{\det g}}
   \partial_i\bigl(\sqrt{\det g}\,g^{ij}\partial_jF\bigr)
\]
therefore gives the first expression in
\eqref{eq:hyperbolic-laplace-beltrami}; the identity
$\partial_x=(\lambda/\mu)\partial_X$ gives the second.  Finally,
contracting $dF$ with the inverse metric and then multiplying by
$d\vol_g$ proves the first integral in
\eqref{eq:hyperbolic-dirichlet-energy}.  Changing from $x$ to $X$ proves the
last one.  This is the two-dimensional conformal invariance of Dirichlet
energy written with all AEG scale factors visible.
\end{proof}

The positively oriented orthonormal arithmetic frame is
\begin{equation}\label{eq:hyperbolic-arithmetic-frame-II}
 \Xu=-\mu y\partial_x=-\lambda y\partial_X,
 \qquad
 \Xv=-\lambda y\partial_y.
\end{equation}
It satisfies
\begin{equation}\label{eq:hyperbolic-frame-bracket-II}
 [\Xu,\Xv]=\lambda\Xu.
\end{equation}
Consequently, the general surface formula of
Section~\ref{sec:surface-operators} becomes
\begin{equation}\label{eq:hyperbolic-laplacian-frame}
 \Delta_g=\Xu^2+\Xv^2+\lambda\Xv.
\end{equation}
The drift $\lambda\Xv$ cancels the first-order term created when $\Xv^2$
is expanded in the $y$-coordinate.  Thus the bare sum
$\Xu^2+\Xv^2$ is not the Laplace--Beltrami operator on this surface.

\subsection{Complex structure and the assignment}

Let
\[
 \partial_{\bar w}=\frac12(\partial_X+i\partial_y),
 \qquad
 \partial_w=\frac12(\partial_X-i\partial_y).
\]
Equations \eqref{eq:hyperbolic-arithmetic-frame-II} give
\begin{equation}\label{eq:hyperbolic-dbar-coordinate}
 \dbarAEG
 =\frac12(\Xu+i\Xv)
 =-\lambda y\,\partial_{\bar w},
 \qquad
 \dAEG=-\lambda y\,\partial_w.
\end{equation}
Since $\lambda y$ never vanishes, arithmetic holomorphicity on the basic
model is exactly ordinary holomorphicity in $w$.  In particular, its real and
imaginary parts are $\Delta_g$-harmonic.

The assignment has the following two useful forms:
\begin{equation}\label{eq:hyperbolic-assignment-w}
 a=-\frac{\mu}{\lambda}\frac{X}{y},
 \qquad
 w=y\left(i-\frac{\lambda}{\mu}a\right).
\end{equation}
It is not itself harmonic.  Directly from
\eqref{eq:hyperbolic-laplace-beltrami},
\begin{equation}\label{eq:hyperbolic-assignment-eigen-equation-II}
 |\nabla a|_g^2=\mu^2+\lambda^2a^2,
 \qquad
 \Delta_ga=2\lambda^2a.
\end{equation}
The second identity is a pointwise eigen-equation in the sign convention of
Paper~I.  It is not an $L^2$ spectral assertion: $a$ is unbounded, has no
finite continuous trace on the ideal boundary, and does not have finite global
Dirichlet energy.

\subsection{The ideal boundary}

The conformal compactification is
\[
 \overline\HH^{\,\mathrm{conf}}
 =\HH\cup\R\cup\{\infty\},
\]
where $\R$ is represented by $y=0$ and the two ends of that line meet at the
single point $\infty$.  This is the closed disc after a Cayley transform.  The
line $y=0$ is not a finite Riemannian boundary component: it lies at infinite
$g_{\mu,\lambda}$-distance.  Boundary conditions below are therefore
conformal or ideal-boundary conditions.

The normalized coordinate $X=(\lambda/\mu)x$ extends to a homeomorphism of
the compactified boundary and fixes the point $\infty$.  Thus the classical
upper-half-plane kernels may be transported to the original arithmetic
$x$-coordinate, but the boundary density must be transformed as well.  That
transformation is carried out explicitly in the next section.
